% self_control.tex
% Питання та тести для самоконтролю з курсу "Основи програмування на Python"
% для студентів першого курсу спеціальності "Математика"
% Оформлення відповідає вимогам ОНУ та стилю файлів main.tex і l1.tex

\chapter{Матеріали для самоконтролю: питання та тести}
\justifying

Цей розділ містить комплекс контрольних запитань та тестових завдань для самостійної перевірки знань і підготовки до заліково-екзаменаційних заходів. Завдання орієнтовані на студентів математичних спеціальностей і спрямовані на глибоке розуміння семантики мови \Python, аналіз інваріантів, обчислювальної складності алгоритмів, особливостей машинної арифметики та формальної логіки.

% ==============================================================================
% РОЗДІЛ 1. ПРОГРАМУВАННЯ ТА АЛГОРИТМИ
% ==============================================================================
\section{Тема 1. Теоретичні основи інформатики та алгоритмізації}

\subsection*{Контрольні запитання}
\begin{enumerate}[label=\textbf{1.\arabic*.}]
	\item Дайте строго-інтуїтивне означення алгоритму. Охарактеризуйте його фундаментальні властивості: дискретність, детермінованість (визначеність), масовість, результативність та скінченність.
	\item Чому властивість детермінованості є критичною при чисельному розв'язанні диференціальних рівнянь та оптимізаційних задач?
	\item Сформулюйте теорему Бема---Якопіні про структурне програмування. Які три базові керуючі структури є функціонально повними для опису будь-якого обчислюваного алгоритму?
	\item Що таке інваріант циклу? Як метод математичної індукції застосовується для доведення коректності та збіжності циклічних алгоритмів?
	\item У чому полягає відмінність між часовою складністю алгоритму $O(f(n))$ та просторовою складністю? Чому при асимптотичному аналізі нехтують константними множниками?
	\item Порівняйте лінійні алгоритми з розгалуженими та циклічними за топологією їхніх блок-схем і графом переходів станів.
\end{enumerate}

\subsection*{Тести для самоконтролю}

\begin{pr}{Тест 1.1. Властивості алгоритмів}
У математичній моделі досліджується ітераційний процес знаходження кореня рівняння $f(x)=0$. Якщо для деяких вхідних даних алгоритм не збігається і виконує нескінченну кількість кроків, яка фундаментальна властивість алгоритму порушується?
\begin{enumerate}[label=\textbf{\Alph*.}]
	\item Дискретність
	\item Детермінованість
	\item Скінченність (результативність)
	\item Масовість
\end{enumerate}
\end{pr}

\begin{pr}{Тест 1.2. Асимптотична складність}
Лінійний алгоритм містить $k$ послідовних операцій введення даних, $m$ елементарних арифметичних дій та виведення результату, де $k, m \in \mathbb{N}$ є фіксованими константами, що не залежать від вхідних величин. Якою є асимптотична часова складність такого алгоритму?
\begin{enumerate}[label=\textbf{\Alph*.}]
	\item $O(1)$
	\item $O(k + m)$
	\item $O(n)$
	\item $O(\log(k \cdot m))$
\end{enumerate}
\end{pr}

\begin{pr}{Тест 1.3. Топологія базових конструкцій}
Яка з наведених блок-схем відповідає принципам структурного програмування за схемою Бема---Якопіні?
\begin{enumerate}[label=\textbf{\Alph*.}]
	\item Конструкція з довільними безумовними переходами (\texttt{goto}) між гілками
	\item Блок, що має рівно один вхід і рівно один вихід
	\item Конструкція з кількома незалежними точками входу в середину циклу
	\item Схема з рекурсивним переходом без термінального вузла
\end{enumerate}
\end{pr}

% ==============================================================================
% РОЗДІЛ 2. ЕКОСИСТЕМА PYTHON ТА ІНТЕРПРЕТАТОР
% ==============================================================================
\section{Тема 2. Архітектура та екосистема Python. Інтерпретатор і середовище}

\subsection*{Контрольні запитання}
\begin{enumerate}[label=\textbf{2.\arabic*.}]
	\item Опишіть етапи трансляції та виконання коду в класичній реалізації CPython: компіляція вихідного тексту у байт-код (\texttt{.pyc}) та його інтерпретація віртуальною машиною PVM (\textit{Python Virtual Machine}).
	\item У чому полягає відмінність між компільованими мовами (C/C++, Fortran) та інтерпретованими (Python) з точки зору швидкодії обчислень та гнучкості розробки?
	\item Що означає характеристика Python як мови із \textit{динамічною}, але \textit{строгою (сильною)} типізацією? Наведіть приклад, що ілюструє суворість типів.
	\item Поясніть концепцію об'єктної моделі пам'яті в Python. Чим відрізняється ідентичність об'єкта (\texttt{id()}) від його числового значення?
	\item Які аргументи має функція \texttt{print()}, і як використання параметрів \texttt{sep} та \texttt{end} дозволяє гнучко форматувати табличний вивід числових даних?
\end{enumerate}

\subsection*{Тести для самоконтролю}

\begin{pr}{Тест 2.1. Форматування виводу функцією \texttt{print()}}
Який результат буде надруковано на екрані після виконання наступного фрагмента коду?
\begin{py}{Лістинг до тесту 2.1}
a, b, c = 1, 2, 3
print(a, b, c, sep="-", end="=")
print(a + b + c)
\end{py}
\begin{enumerate}[label=\textbf{\Alph*.}]
	\item \texttt{1-2-3=6}
	\item \texttt{1-2-3} \\ \texttt{=6}
	\item \texttt{1 2 3 = 6}
	\item \texttt{1-2-3= 6}
\end{enumerate}
\end{pr}

\begin{pr}{Тест 2.2. Строга динамічна типізація}
Яким буде наслідок спроби обчислення виразу \texttt{"10" + 20} у консолі Python без попереднього перетворення типів?
\begin{enumerate}[label=\textbf{\Alph*.}]
	\item Рядок \texttt{"1020"} (неявне приведення до рядка)
	\item Число \texttt{30} (неявне приведення до числа)
	\item Виникнення виключення \texttt{TypeError}
	\item Число \texttt{200}
\end{enumerate}
\end{pr}

\begin{pr}{Тест 2.3. Посилальна модель пам'яті}
Що буде виведено після виконання послідовності інструкцій?
\begin{py}{Лістинг до тесту 2.3}
x = 500
y = x
x = x + 1
print(x, y, x is y)
\end{py}
\begin{enumerate}[label=\textbf{\Alph*.}]
	\item \texttt{501 501 True}
	\item \texttt{501 500 False}
	\item \texttt{501 500 True}
	\item \texttt{501 501 False}
\end{enumerate}
\end{pr}

% ==============================================================================
% РОЗДІЛ 3. ТИПИ ДАНИХ ТА АРИФМЕТИКА
% ==============================================================================
\section{Тема 3. Синтаксис Python. Базові типи даних та арифметика}

\subsection*{Контрольні запитання}
\begin{enumerate}[label=\textbf{3.\arabic*.}]
	\item Як реалізовано довгу цілочисельну арифметику (\texttt{int}) у Python 3? Чому обчислення факторіала $100!$ не викликає переповнення розрядної сітки?
	\item Сформулюйте причини похибок при роботі з типом \texttt{float} (стандарт IEEE 754). Чому рівність \texttt{0.1 + 0.2 == 0.3} є хибною, і як коректно перевіряти рівність дійсних чисел $x, y \in \mathbb{R}$ з точністю $\varepsilon$?
	\item Сформулюйте правило цілочисельного ділення \texttt{//} та взяття остачі \texttt{\%} у Python. Як воно узгоджується з фундаментальною математичною формулою $a = b \cdot q + r$, де частка $q = \lfloor a/b \rfloor$, а для залишку гарантується $0 \le r < b$ (для $b > 0$)?
	\item Обчисліть аналітично та за правилами Python значення виразів \texttt{-17 // 5} та \texttt{-17 \% 5}. Чому результат відрізняється від мов C/C++?
	\item Поясніть концепцію «банківського округлення» (\textit{round half to even}), реалізованого у вбудованій функції \texttt{round()}. Порівняйте значення \texttt{round(2.5)} та \texttt{round(3.5)}.
	\item Опишіть алгоритм виділення розрядів довільного трицифрового натурального числа $N \in [100, 999]$ за допомогою комбінації операцій \texttt{//} та \texttt{\%}.
\end{enumerate}

\subsection*{Тести для самоконтролю}

\begin{pr}{Тест 3.1. Цілочисельне ділення від'ємних чисел}
Які значення поверне Python для наступного виклику:
\begin{py}{Лістинг до тесту 3.1}
print(-23 // 4, -23 % 4)
\end{py}
\begin{enumerate}[label=\textbf{\Alph*.}]
	\item \texttt{-5 -3}
	\item \texttt{-6 1}
	\item \texttt{-6 -3}
	\item \texttt{-5 1}
\end{enumerate}
\end{pr}

\begin{pr}{Тест 3.2. Правоасоціативність оператора піднесення до степеня}
Чому дорівнює значення математичного виразу в Python:
\begin{py}{Лістинг до тесту 3.2}
val = 2 ** 3 ** 2
print(val)
\end{py}
\begin{enumerate}[label=\textbf{\Alph*.}]
	\item $64$ (оскільки $(2^3)^2 = 8^2 = 64$)
	\item $512$ (оскільки $2^{(3^2)} = 2^9 = 512$)
	\item $36$
	\item $128$
\end{enumerate}
\end{pr}

\begin{pr}{Тест 3.3. Банківське округлення чисел}
Який кортеж значень утворить наступний вираз:
\begin{py}{Лістинг до тесту 3.3}
res = (round(4.5), round(5.5), round(-2.5))
print(res)
\end{py}
\begin{enumerate}[label=\textbf{\Alph*.}]
	\item \texttt{(5, 6, -3)}
	\item \texttt{(4, 6, -2)}
	\item \texttt{(5, 5, -2)}
	\item \texttt{(4, 5, -3)}
\end{enumerate}
\end{pr}

\begin{pr}{Тест 3.4. Виділення десяткових розрядів}
Дано число $n = 783$. Який вираз поверне середню цифру числа (тобто цифру 8)?
\begin{enumerate}[label=\textbf{\Alph*.}]
	\item \texttt{n // 100}
	\item \texttt{n \% 10}
	\item \texttt{(n // 10) \% 10}
	\item \texttt{n // (10 \% 10)}
\end{enumerate}
\end{pr}

% ==============================================================================
% РОЗДІЛ 4. УМОВНИЙ ОПЕРАТОР IF
% ==============================================================================
\section{Тема 4. Розгалужені алгоритми: умовний оператор if}

\subsection*{Контрольні запитання}
\begin{enumerate}[label=\textbf{4.\arabic*.}]
	\item Які існують синтаксичні форми умовного оператора в Python (неповне, повне та багатоальтернативне розгалуження \texttt{if-elif-else})?
	\item Чому в конструкції \texttt{if-elif-else} порядок запису умов є істотним? Наведіть приклад, коли неправильний порядок умов призводить до недосяжності окремих гілок коду.
	\item Як реалізовано ланцюгові порівняння в Python виду $a < x \le b$? Скільки разів обчислюється значення змінної чи виразу $x$ при такій перевірці?
	\item Запишіть у термінах мови Python розгалужений алгоритм обчислення кусочно-неперервної функції:
	\[
	f(x) = \begin{cases}
		\sin(x), & x < 0, \\
		x^2 + 1, & 0 \le x \le 2, \\
		\sqrt{x}, & x > 2.
	\end{cases}
	\]
	\item Як працює тернарний умовний оператор виразу \texttt{x = A if cond else B}? За яких умов його застосування покращує читабельність коду?
\end{enumerate}

\subsection*{Тести для самоконтролю}

\begin{pr}{Тест 4.1. Ланцюгові порівняння}
Що надрукує інтерпретатор при виконанні наступного виразу?
\begin{py}{Лістинг до тесту 4.1}
x = 5
if 1 < x <= 5 == 5:
    print("ІСТИНА")
else:
    print("ХИБА")
\end{py}
\begin{enumerate}[label=\textbf{\Alph*.}]
	\item \texttt{ІСТИНА}
	\item \texttt{ХИБА}
	\item Виникне синтаксична помилка \texttt{SyntaxError}
	\item Виникне помилка типів \texttt{TypeError}
\end{enumerate}
\end{pr}

\begin{pr}{Тест 4.2. Порядок гілок в \texttt{if-elif-else}}
Нехай задано ціле число $x = 15$. Що виведе наступна програма?
\begin{py}{Лістинг до тесту 4.2}
x = 15
if x > 5:
    print("A", end="")
elif x > 10:
    print("B", end="")
else:
    print("C", end="")
\end{py}
\begin{enumerate}[label=\textbf{\Alph*.}]
	\item \texttt{AB}
	\item \texttt{B}
	\item \texttt{A}
	\item \texttt{ABC}
\end{enumerate}
\end{pr}

\begin{pr}{Тест 4.3. Тернарний вираз та присвоєння}
Якого значення набуде змінна $y$ в результаті виконання коду?
\begin{py}{Лістинг до тесту 4.3}
a, b = 7, 12
y = a * 2 if a > b else b // 3
print(y)
\end{py}
\begin{enumerate}[label=\textbf{\Alph*.}]
	\item $14$
	\item $4$
	\item $0$
	\item $True$
\end{enumerate}
\end{pr}

% ==============================================================================
% РОЗДІЛ 5. ЛОГІЧНІ ОПЕРАТОРИ
% ==============================================================================
\section{Тема 5. Логічні оператори та булева алгебра}

\subsection*{Контрольні запитання}
\begin{enumerate}[label=\textbf{5.\arabic*.}]
	\item Сформулюйте таблиці істинності для логічних зв'язок \texttt{and} (кон'юнкція), \texttt{or} (диз'юнкція) та \texttt{not} (інверсія).
	\item Запишіть закони де Моргана засобами мови Python. Як за їх допомогою спростити предикат \texttt{not (x >= 0 and x <= 100)}?
	\item У чому полягає принцип «лінивих обчислень» (\textit{short-circuit evaluation}) для операторів \texttt{and} та \texttt{or}? Як цей механізм захищає програму від ділення на нуль: \texttt{if x != 0 and 1/x > 1}?
	\item Які типи об'єктів у Python мають булевий еквівалент \texttt{False} (\textit{Falsy}), а які --- \texttt{True} (\textit{Truthy})? Що повертає виклик \texttt{bool("")}, \texttt{bool(0.0)}, \texttt{bool([])}?
	\item Чому логічні оператори \texttt{and} та \texttt{or} у Python повертають значення одного зі своїх операндів, а не суто тип \texttt{bool}? Що є результатом виразу \texttt{"alpha" or "beta"}?
\end{enumerate}

\subsection*{Тести для самоконтролю}

\begin{pr}{Тест 5.1. Повернення операндів у лінивих обчисленнях}
Який точний результат буде отримано під час обчислення наступного виразу?
\begin{py}{Лістинг до тесту 5.1}
res = 0 or ("Python" and 42)
print(res)
\end{py}
\begin{enumerate}[label=\textbf{\Alph*.}]
	\item \texttt{True}
	\item \texttt{0}
	\item \texttt{"Python"}
	\item \texttt{42}
\end{enumerate}
\end{pr}

\begin{pr}{Тест 5.2. Пріоритет булевих операцій}
Яке булеве значення поверне наступний логічний вираз?
\begin{py}{Лістинг до тесту 5.2}
val = not True or False and not False
print(val)
\end{py}
\begin{enumerate}[label=\textbf{\Alph*.}]
	\item \texttt{True}
	\item \texttt{False}
	\item \texttt{None}
	\item Виникне помилка пріоритету
\end{enumerate}
\end{pr}

\begin{pr}{Тест 5.3. Захист від аварійного завершення (Short-circuiting)}
Нехай $x = 0$. Який із наведених нижче виразів \textbf{не викличе} помилку \texttt{ZeroDivisionError}?
\begin{enumerate}[label=\textbf{\Alph*.}]
	\item \texttt{10 / x > 1 and x != 0}
	\item \texttt{x != 0 and 10 / x > 1}
	\item \texttt{x == 0 or 10 / x > 1 and False}
	\item Як вираз B, так і вираз C
\end{enumerate}
\end{pr}

% ==============================================================================
% РОЗДІЛ 6. ГЕНЕРАЦІЯ ВИПАДКОВИХ ЧИСЕЛ: МОДУЛЬ RANDOM
% ==============================================================================
\section{Тема 6. Генерація псевдовипадкових чисел: модуль random}

\subsection*{Контрольні запитання}
\begin{enumerate}[label=\textbf{6.\arabic*.}]
	\item Що таке псевдовипадкові числа? На якому базовому алгоритмі (наприклад, \textit{Mersenne Twister}) базується генератор модуля \texttt{random} у Python?
	\item Яку функцію виконує метод \texttt{random.seed(k)}? Чому фіксація зерна генератора є обов'язковою для наукової верифікації та повторюваності чисельних результатів?
	\item У чому різниця між функціями \texttt{random.random()} та \texttt{random.uniform(a, b)}? Які розподіли ймовірностей вони моделюють?
	\item Порівняйте поведінку функцій \texttt{random.randint(a, b)} та \texttt{random.randrange(start, stop, step)}. Чи включаються граничні точки $a$ та $b$ у діапазон можливих результатів?
	\item Опишіть математичну схему методу Монте-Карло для обчислення числа $\pi$ через відношення кількості точок, що потрапили у чверть одиничного круга, до загальної кількості згенерованих точок у квадраті $[0, 1] \times [0, 1]$.
\end{enumerate}

\subsection*{Тести для самоконтролю}

\begin{pr}{Тест 6.1. Межі функцій \texttt{randint} та \texttt{randrange}}
Який з наведених викликів генерує випадкове ціле число $k \in \{1, 2, 3, 4, 5, 6\}$ (моделювання кидка грального кубика)?
\begin{enumerate}[label=\textbf{\Alph*.}]
	\item \texttt{random.randint(1, 6)}
	\item \texttt{random.randrange(1, 6)}
	\item \texttt{random.uniform(1, 6)}
	\item \texttt{int(random.random() * 6)}
\end{enumerate}
\end{pr}

\begin{pr}{Тест 6.2. Фіксація стану генератора}
Що можна стверджувати про виконання наступного коду на двох різних комп'ютерах з однаковою версією Python?
\begin{py}{Лістинг до тесту 6.2}
import random
random.seed(12345)
a = [random.randint(1, 100) for _ in range(5)]
print(a)
\end{py}
\begin{enumerate}[label=\textbf{\Alph*.}]
	\item Списки будуть різними, оскільки апаратний генератор системного часу унікальний
	\item Списки будуть абсолютно ідентичними, оскільки початковий стан псевдовипадкової послідовності зафіксовано
	\item Програма видасть помилку під час повторного запуску
	\item Списки збігатимуться лише якщо довжина списку дорівнює 1
\end{enumerate}
\end{pr}

\begin{pr}{Тест 6.3. Моделювання рівномірного розподілу}
Нехай $\xi = \texttt{random.random()}$. Яка лінійна трансформація відображає $\xi \in [0, 1)$ на неперервний інтервал $[a, b)$ для довільних $a < b$?
\begin{enumerate}[label=\textbf{\Alph*.}]
	\item $a + \xi \cdot b$
	\item $a + \xi \cdot (b - a)$
	\item $\xi \cdot (b + a)$
	\item $b - \xi \cdot a$
\end{enumerate}
\end{pr}

% ==============================================================================
% РОЗДІЛ 7. ЦИКЛ FOR ТА ФУНКЦІЯ RANGE()
% ==============================================================================
\section{Тема 7. Циклічні алгоритми: цикл for та функція range()}

\subsection*{Контрольні запитання}
\begin{enumerate}[label=\textbf{7.\arabic*.}]
	\item Як інтерпретується цикл \texttt{for} у мові Python? Чому його визначають як ітератор по довільній послідовності об'єктів, а не просто як лічильник?
	\item Опишіть три форми виклику функції \texttt{range()}: з одним, двома та трьома параметрами. Чому функція \texttt{range} у Python 3 не займає пам'ять пропорційно до кількості елементів?
	\item Як за допомогою параметрів \texttt{range(start, stop, step)} згенерувати послідовність цілих чисел від 20 до 0 включно з кроком $-2$?
	\item Опишіть шаблон алгоритму обчислення суми ряду $\sum_{k=1}^n a_k$ та добутку $\prod_{k=1}^n a_k$. Якими повинні бути початкові значення змінних-накопичувачів перед входом у цикл?
	\item Як будуються вкладені цикли \texttt{for}? Якою є загальна кількість ітерацій внутрішнього тіла при виклику:
	\begin{pydef}
for i in range(N):
    for j in range(M):
        оператори\end{pydef}
\end{enumerate}

\subsection*{Тести для самоконтролю}

\begin{pr}{Тест 7.1. Кількість ітерацій \texttt{range()}}
Скільки разів виконається тіло наступного циклу?
\begin{py}{Лістинг до тесту 7.1}
count = 0
for k in range(3, 23, 4):
    count += 1
print(count)
\end{py}
\begin{enumerate}[label=\textbf{\Alph*.}]
	\item $4$ рази
	\item $5$ разів (значення $k$: 3, 7, 11, 15, 19)
	\item $6$ разів
	\item $20$ разів
\end{enumerate}
\end{pr}

\begin{pr}{Тест 7.2. Порожній числовий діапазон}
Яким буде виведення програми після запуску наступного коду?
\begin{py}{Лістинг до тесту 7.2}
s = 100
for i in range(10, 2):
    s += i
print(s)
\end{py}
\begin{enumerate}[label=\textbf{\Alph*.}]
	\item $144$
	\item $100$ (тіло циклу не виконається жодного разу через додатний крок за замовчуванням)
	\item $108$
	\item Виникне виключення \texttt{ValueError}
\end{enumerate}
\end{pr}

\begin{pr}{Тест 7.3. Вкладені цикли та акумуляція суми}
Який результат надрукує цей фрагмент програми?
\begin{py}{Лістинг до тесту 7.3}
total = 0
for i in range(1, 4):
    for j in range(i):
        total += 1
print(total)
\end{py}
\begin{enumerate}[label=\textbf{\Alph*.}]
	\item $12$
	\item $9$
	\item $6$ (обчислення суми $0 + 1 + 2 = 3$ та $1+2+3 = 6$)
	\item $3$
\end{enumerate}
\end{pr}

% ==============================================================================
% РОЗДІЛ 8. ЦИКЛ WHILE ТА КЕРУВАННЯ ЦИКЛАМИ
% ==============================================================================
\section{Тема 8. Цикли з передумовою while та керування ітераціями}

\subsection*{Контрольні запитання}
\begin{enumerate}[label=\textbf{8.\arabic*.}]
	\item Сформулюйте семантику оператора циклу \texttt{while}. Які умови мають бути виконані для гарантії завершення циклічного процесу?
	\item Поясніть призначення та дію інструкцій керування циклом: \texttt{break}, \texttt{continue} та \texttt{pass}.
	\item Як працює конструкція \texttt{while ... else} у Python? За якої умови гілка \texttt{else} \textbf{не} виконується?
	\item Опишіть алгоритм Евкліда для знаходження найбільшого спільного дільника $\gcd(a, b)$ через операцію остачі від ділення. Якою є його логарифмічна складність?
	\item Чому для перевірки простоти натурального числа $N$ достатньо шукати дільники лише в діапазоні $2 \le d \le \lfloor\sqrt{N}\rfloor$?
	\item Чому при чисельному обчисленні границь послідовностей чи рядів небезпечно писати умову \texttt{while x != target:}, якщо \texttt{x} має тип \texttt{float}?
\end{enumerate}

\subsection*{Тести для самоконтролю}

\begin{pr}{Тест 8.1. Конструкція \texttt{while ... else}}
Що буде надруковано на екрані в результаті виконання наступного коду?
\begin{py}{Лістинг до тесту 8.1}
n = 5
while n > 0:
    n -= 2
    if n == 1:
        break
else:
    print("DONE", end=" ")
print("END")
\end{py}
\begin{enumerate}[label=\textbf{\Alph*.}]
	\item \texttt{DONE END}
	\item \texttt{END} (блок \texttt{else} пропускається, бо цикл перервано через \texttt{break})
	\item \texttt{DONE}
	\item Нескінченний цикл
\end{enumerate}
\end{pr}

\begin{pr}{Тест 8.2. Дія оператора \texttt{continue}}
Чому дорівнюватиме значення змінної $s$ після завершення циклу?
\begin{py}{Лістинг до тесту 8.2}
s, i = 0, 0
while i < 6:
    i += 1
    if i % 2 == 0:
        continue
    s += i
print(s)
\end{py}
\begin{enumerate}[label=\textbf{\Alph*.}]
	\item $21$
	\item $9$ (додаються лише непарні числа: $1 + 3 + 5 = 9$)
	\item $12$
	\item $6$
\end{enumerate}
\end{pr}

\begin{pr}{Тест 8.3. Кількість операцій в алгоритмі Евкліда}
Скільки ітерацій циклу \texttt{while} знадобиться алгоритму Евкліда для знаходження $\gcd(48, 18)$?
\begin{py}{Лістинг до тесту 8.3}
a, b = 48, 18
k = 0
while b != 0:
    a, b = b, a % b
    k += 1
print(k, a)
\end{py}
\begin{enumerate}[label=\textbf{\Alph*.}]
	\item $3$ ітерації, результат $\gcd = 6$
	\item $2$ ітерації, результат $\gcd = 6$
	\item $4$ ітерації, результат $\gcd = 2$
	\item $1$ ітерація, результат $\gcd = 18$
\end{enumerate}
\end{pr}

% ==============================================================================
% РОЗДІЛ 9. СПИСКИ (LIST) ТА АЛГОРИТМИ ЇХ ОБРОБКИ
% ==============================================================================
\section{Тема 9. Базові структури даних: списки (list) та генератори списків}

\subsection*{Контрольні запитання}
\begin{enumerate}[label=\textbf{9.\arabic*.}]
	\item Охарактеризуйте тип даних \texttt{list} у Python. Які його фундаментальні відмінності від статичних масивів у C/C++?
	\item Як влаштовані зрізи списків \texttt{L[start:stop:step]}? Що повертає інструкція \texttt{L[::-1]}? Чи створює операція зрізу копію елементів?
	\item Поясніть концепцію генераторів списків (\textit{list comprehension}). Запишіть конструктор для створення списку квадратів непарних чисел у діапазоні $[1, 50]$.
	\item Чому створення двовимірного списку (матриці) розміром $n \times m$ через вираз \texttt{M = [[0] * m] * n} є грубою помилкою? Як правильно ініціалізувати нульову матрицю?
	\item Яка асимптотична складність операцій додавання елемента в кінець списку (\texttt{.append()}) та вставки на початок списку (\texttt{.insert(0, val)})?
\end{enumerate}

\subsection*{Тести для самоконтролю}

\begin{pr}{Тест 9.1. Пастка розмноження посилань у матрицях}
Що надрукує наступний фрагмент коду?
\begin{py}{Лістинг до тесту 9.1}
matrix = [[0] * 3] * 3
matrix[0][1] = 9
print(matrix[2][1])
\end{py}
\begin{enumerate}[label=\textbf{\Alph*.}]
	\item $0$
	\item $9$ (усі рядки посилаються на один і той самий об'єкт у пам'яті)
	\item Виникне виключення \texttt{IndexError}
	\item \texttt{None}
\end{enumerate}
\end{pr}

\begin{pr}{Тест 9.2. Зрізи з від'ємним кроком}
Нехай $L = [10, 20, 30, 40, 50, 60]$. Чому дорівнює значення виразу \texttt{L[4:1:-2]}?
\begin{enumerate}[label=\textbf{\Alph*.}]
	\item \texttt{[50, 30]}
	\item \texttt{[50, 40, 30]}
	\item \texttt{[40, 20]}
	\item \texttt{[]}
\end{enumerate}
\end{pr}

\begin{pr}{Тест 9.3. Генератор списку з умовою}
Який список буде сформовано в результаті виконання виразу:
\begin{py}{Лістинг до тесту 9.3}
res = [x // 2 for x in range(10) if x % 3 == 0]
print(res)
\end{py}
\begin{enumerate}[label=\textbf{\Alph*.}]
	\item \texttt{[0, 1, 3, 4]} (числа 0, 3, 6, 9 діляться націло на 2)
	\item \texttt{[0, 1.5, 3.0, 4.5]}
	\item \texttt{[0, 3, 6, 9]}
	\item \texttt{[1, 2, 4]}
\end{enumerate}
\end{pr}

% ==============================================================================
% РОЗДІЛ 10. ФУНКЦІЇ ТА РЕКУРСІЯ
% ==============================================================================
\section{Тема 10. Функції у Python: сигнатура, параметри та рекурсія}

\subsection*{Контрольні запитання}
\begin{enumerate}[label=\textbf{10.\arabic*.}]
	\item У чому полягає відмінність між формальними та фактичними параметрами функції? Як працює механізм передачі аргументів за об'єктним посиланням (\textit{call by object reference})?
	\item Чому використання змінюваних об'єктів (наприклад, \texttt{def f(lst=[])}) як значень параметрів за замовчуванням є типовою пасткою в Python?
	\item Сформулюйте математичні засади рекурсії. Як рекурсивне визначення функції пов'язане з принципом математичної індукції?
	\item Що таке глибина рекурсії та стек викликів? Чому наївна рекурсивна реалізація чисел Фібоначчі $F_n = F_{n-1} + F_{n-2}$ має експоненційну складність $O(2^n)$?
	\item Опишіть правила пошуку імен за правилом LEGB (Local, Enclosing, Global, Built-in).
\end{enumerate}

\subsection*{Тести для самоконтролю}

\begin{pr}{Тест 10.1. Змінювані параметри за замовчуванням}
Що буде виведено після послідовних викликів функції?
\begin{py}{Лістинг до тесту 10.1}
def add_item(val, acc=[]):
    acc.append(val)
    return acc

print(add_item(1))
print(add_item(2))
\end{py}
\begin{enumerate}[label=\textbf{\Alph*.}]
	\item \texttt{[1]} \\ \texttt{[2]}
	\item \texttt{[1]} \\ \texttt{[1, 2]} (значення за замовчуванням створюється один раз під час компіляції функції)
	\item \texttt{[1, 2]} \\ \texttt{[1, 2]}
	\item Виникне виключення \texttt{RuntimeError}
\end{enumerate}
\end{pr}

\begin{pr}{Тест 10.2. Рекурсивне обчислення}
Яке значення поверне виклик функції \texttt{calc(4)}?
\begin{py}{Лістинг до тесту 10.2}
def calc(n):
    if n <= 1:
        return 1
    return n * calc(n - 2)

print(calc(5))
\end{py}
\begin{enumerate}[label=\textbf{\Alph*.}]
	\item $120$ ($5!$)
	\item $15$ (подвійний факторіал $5!! = 5 \cdot 3 \cdot 1$)
	\item $24$
	\item $1$
\end{enumerate}
\end{pr}

\begin{pr}{Тест 10.3. Область видимості змінних}
Що надрукує програма після виклику функції?
\begin{py}{Лістинг до тесту 10.3}
x = 10
def modify():
    x = 20
modify()
print(x)
\end{py}
\begin{enumerate}[label=\textbf{\Alph*.}]
	\item $20$
	\item $10$ (змінна всередині функції є локальною і не модифікує глобальну)
	\item \texttt{None}
	\item Виникне помилка \texttt{UnboundLocalError}
\end{enumerate}
\end{pr}

% ==============================================================================
% РОЗДІЛ: ВІДПОВІДІ ТА АНАЛІЗ ТЕСТІВ
% ==============================================================================
\section{Відповіді та докладний аналіз тестових завдань}

\renewcommand{\arraystretch}{1.2}
\begin{longtable}{|c|c|p{11.5cm}|}
	\hline
	\textbf{Тест} & \textbf{Відповідь} & \textbf{Коротке математичне та алгоритмічне обґрунтування} \\ \hline
	\endfirsthead
	\hline
	\textbf{Тест} & \textbf{Відповідь} & \textbf{Коротке математичне та алгоритмічне обґрунтування} \\ \hline
	\endhead
	\hline
	\endfoot
	\textbf{1.1} & \textbf{C} & Нескінченне виконання порушує властивість результативності (скінченності за фіксований час). \\ \hline
	\textbf{1.2} & \textbf{A} & Кількість кроків фіксована і не залежить від розміру вхідних даних: $T(n) = C \Rightarrow O(1)$. \\ \hline
	\textbf{1.3} & \textbf{B} & За теоремою Бема---Якопіні структурні блоки повинні мати строго один вхід і один вихід. \\ \hline
	\textbf{2.1} & \textbf{A} & \texttt{sep="-"} розмежовує елементи дефісом, а \texttt{end="="} замінює перехід рядка на знак рівності. \\ \hline
	\textbf{2.2} & \textbf{C} & Python має сувору типізацію: додавання рядка до числа не має неявної конвертації типів. \\ \hline
	\textbf{2.3} & \textbf{B} & Цілі числа незмінні (\textit{immutable}). Вираз \texttt{x = x + 1} створює новий об'єкт в пам'яті. \\ \hline
	\textbf{3.1} & \textbf{B} & За правилом округлення вниз: $\lfloor -23/4 \rfloor = -6$. Остача: $-23 - 4 \cdot (-6) = 1$. \\ \hline
	\textbf{3.2} & \textbf{B} & Операція \texttt{**} правоасоціативна: $2^{3^2} = 2^9 = 512$. \\ \hline
	\textbf{3.3} & \textbf{B} & Банківське округлення округлює до найближчого парного: $4.5 \to 4$, $5.5 \to 6$, $-2.5 \to -2$. \\ \hline
	\textbf{3.4} & \textbf{C} & \texttt{783 // 10 = 78}, далі \texttt{78 \% 10 = 8}. \\ \hline
	\textbf{4.1} & \textbf{A} & Ланцюгове порівняння еквівалентне: $(1 < 5) \land (5 \le 5) \land (5 == 5) \Rightarrow \text{True}$. \\ \hline
	\textbf{4.2} & \textbf{C} & Перша істинна гілка $x > 5$ спрацьовує і завершує оператор; блок $x > 10$ ігнорується. \\ \hline
	\textbf{4.3} & \textbf{B} & Оскільки $7 > 12$ є $\text{False}$, обчислюється гілка \texttt{else}: $12 // 3 = 4$. \\ \hline
	\textbf{5.1} & \textbf{D} & У виразі \texttt{"Python" and 42} лівий операнд істинний, результат 42. Далі \texttt{0 or 42} дає 42. \\ \hline
	\textbf{5.2} & \textbf{B} & Пріоритет: $\text{not True} \to \text{False}$; $\text{not False} \to \text{True}$; $\text{False and True} \to \text{False}$; $\text{False or False} \to \text{False}$. \\ \hline
	\textbf{5.3} & \textbf{D} & Обидва вирази B і C захищені лінивими обчисленнями: ділення на нуль не викликається. \\ \hline
	\textbf{6.1} & \textbf{A} & \texttt{randint(a, b)} генерує значення на замкненому відрізку $[a, b]$, включаючи 6. \\ \hline
	\textbf{6.2} & \textbf{B} & Алгоритм Mersenne Twister є повністю детермінованим при фіксації значення seed. \\ \hline
	\textbf{6.3} & \textbf{B} & Формула лінійного масштабування рівномірного розподілу з $[0, 1)$ на $[a, b)$. \\ \hline
	\textbf{7.1} & \textbf{B} & Значення лічильника: 3, 7, 11, 15, 19 (рівно 5 значень, бо $19 + 4 = 23 \ge 23$). \\ \hline
	\textbf{7.2} & \textbf{B} & За замовчуванням крок $+1$. Оскільки $10 > 2$, діапазон порожній, цикл не виконується. \\ \hline
	\textbf{7.3} & \textbf{C} & Внутрішній цикл виконується $1 + 2 + 3 = 6$ разів. \\ \hline
	\textbf{8.1} & \textbf{B} & Блок \texttt{else} не виконується, якщо цикл завершено аварійно через команду \texttt{break}. \\ \hline
	\textbf{8.2} & \textbf{B} & \texttt{continue} пропускає додавання парних чисел; сума непарних: $1 + 3 + 5 = 9$. \\ \hline
	\textbf{8.3} & \textbf{A} & Кроки: $(48, 18) \to (18, 12) \to (12, 6) \to (6, 0)$. Рівно 3 кроки, $\gcd = 6$. \\ \hline
	\textbf{9.1} & \textbf{B} & Множення списку створює не копії, а однакові посилання на єдиний підсписок у пам'яті. \\ \hline
	\textbf{9.2} & \textbf{A} & Індекс 4 (50), крок $-2 \implies$ індекс 2 (30). Верхня межа 1 не включається. \\ \hline
	\textbf{9.3} & \textbf{A} & Відібрані числа $x \in \{0, 3, 6, 9\}$. Цілочисельне ділення на 2 дає \texttt{[0, 1, 3, 4]}. \\ \hline
	\textbf{10.1} & \textbf{B} & Об'єкт списку за замовчуванням зв'язується один раз і накопичує елементи між викликами. \\ \hline
	\textbf{10.2} & \textbf{B} & Обчислюється подвійний факторіал непарного числа: $5 \cdot 3 \cdot 1 = 15$. \\ \hline
	\textbf{10.3} & \textbf{B} & Присвоєння $x = 20$ створило локальну змінну всередині функції; глобальне $x$ лишилося 10. \\ \hline
\end{longtable}
